\input{preamble}

\begin{document}

\title{From Obstruction to Rate:\\The Observation Order of a Safe Compression Decision}

\author{\IEEEauthorblockN{Vu Phan An Nguyen}
\IEEEauthorblockA{\textit{Rider University} \\
Lawrenceville, NJ, USA}}

\maketitle

\begin{abstract}
A learner that must choose a compression does not have to identify the world it is in. It
only has to separate worlds far enough that no surviving set of them disagrees on every
answer. We turn that into a rate law and prove it.

Call a set of worlds with no common safe answer an \emph{obstruction}. For a finite world
set, take any safe-answer correspondence and any channel whose queries carry at most
$\kappa(w,w')$ nats between two worlds. The cost of a safe decision is then governed by one
number: $t^{*}$, the minimum over obstructions of the largest information among their pairs.
Every finite-horizon policy that answers safely with probability $1-\delta$ spends, at some
world, at least $\bigl(\frac{1-\delta}{|W|-1}\log\frac1\delta-\log2\bigr)/t^{*}$ queries in
expectation. A tournament that tests every pair of information at least $t^{*}$ attains
$O(|W|^{2}\log(|W|/\delta)/t^{*})$ whenever pairwise tests have exponent comparable to
$\kappa$.

On a graded channel, with $\kappa\asymp h^{2\,\mathrm{sep}}$, this becomes an \emph{order law}:
$t^{*}=Ch^{2s^{*}}$. Here $s^{*}$ is the least order at which every set of worlds the observer
still confuses has a common safe answer, that is, the maximum over obstructions of their
minimum separation order. Identification is the special case of singleton answer sets, and its
order is the maximum separation.

On the parity-mask channel with $m$ worlds both rates are theorems, against every randomised
adaptive policy that stops almost surely. The safe decision costs
$\Theta(\log(1/\delta)/h^{2})$ and identification costs $\Theta(\log(1/\delta)/h^{2(m-1)})$. Both
lower sides are corollaries of the law. The boundary is a second family with a latent world,
where the per-step information orders $h^{2}$ and $h^{2r}$ are proved uniformly over the linear
query class, and the query class itself decides the order. All theorems are machine-checked in
Lean~4.
\end{abstract}

\begin{IEEEkeywords}
sequential design, adaptive lower bounds, decision regions, information orders, formal
verification
\end{IEEEkeywords}

\section{Introduction}
\label{sec:intro}

\subsection{The question}

Identifying the world suffices for choosing an action, and it is usually far more than
necessary. The question is how much more, and whether the gap can be predicted from the
structure of the problem rather than measured after the fact.

Fix $\varepsilon$, and let $\mathcal A(w)$ be the set of $\varepsilon$-safe compressions in
world $w$. A set of worlds $H$ is an \emph{obstruction} when
$\bigcap_{w\in H}\mathcal A(w)=\varnothing$. A learner that cannot yet tell apart the worlds of
an obstruction cannot answer safely, whatever it randomises. This paper proves the law that
turns obstructions and the information of the queries into a rate. It also proves the law's
two-rate instance on one channel, and it maps the boundary where the query class, rather than
the worlds, decides the order.

\subsection{Results}

\begin{enumerate}
\item[1)] \textbf{The obstruction law} (Theorems~\ref{thm:lower}, \ref{thm:upper}). For every
finite world set, safe-answer correspondence and channel, the minimax cost is governed by
$t^{*}$, the minimum over obstructions of the largest per-query information among their pairs.
The lower bound $\Omega(\log(1/\delta)/t^{*})$ holds for every channel and every policy. A
tournament attains $O(\log(|W|/\delta)/t^{*})$ given pairwise tests of comparable exponent.
\item[2)] \textbf{The order law} (Theorem~\ref{thm:order}). On a graded channel $t^{*}=Ch^{2s^{*}}$,
with $s^{*}$ the least safe order. For overlapping safe sets, $s^{*}$ is the minimum
separation order. For identification it is the maximum.
\item[3)] \textbf{Both rates on one channel} (Theorems~\ref{thm:endpoint},
\ref{thm:parlaw}). On the parity-mask channel: $\Theta(\log(1/\delta)/h^{2})$ for the safe
decision and $\Theta(\log(1/\delta)/h^{2(m-1)})$ for identification, against every randomised
adaptive policy that stops almost surely. The lower sides follow from the law.
\item[4)] \textbf{The query class decides the order} (Theorems~\ref{thm:vanish},
\ref{thm:attain}, \ref{thm:unif}). On a latent family, the signal against the hidden character
vanishes below order $\lceil n/d\rceil$ for every query of degree $d$, and the averaged tail
query attains that order with constant exactly $r!/r^{r}$. The per-step information is at most
$Ch^{2r}$ uniformly over the linear class. A degree-capped class breaks it, and one parameter
of nonlinearity moves the order from $h^{6}$ to $h^{4}$.
\end{enumerate}

\subsection{Prior work, and what does not follow from it}

Deciding without identifying is not ours. Javdani et al.~\cite{javdani2014} formulate decisions
with overlapping regions and give near-optimal Bayesian active learning for them. Signal that
appears only at high moments is the masking mechanism of Grosso, Standaert and
Prouff~\cite{grosso2014}, with moment expansions of leakage by Rioul, Cheng and
Guilley~\cite{rioul2021}, and concrete masking bounds by Duc, Faust and Standaert~\cite{duc2015}.
Keeping only decision-relevant information is itself established~\cite{poli2023}. Information
games with several correct answers are developed by Degenne and Koolen~\cite{degenne2019}, and
ranking-and-selection with multiple correct answers and non-answerable estimates by Wang and
You~\cite{wangyou2026}. We use these as tools and do not claim them.

What is ours is the law that joins them. Obstructions and per-query information determine the
rate for every safe-answer correspondence and every channel, and a graded channel turns that
into an order.

\begin{table}[h]
\caption{Nearest sources, and a theorem of this paper that does not follow from each.}
\label{tab:prior}
\centering
\footnotesize
\begin{tabular}{@{}p{0.30\columnwidth}p{0.62\columnwidth}@{}}
\toprule
Nearest source & Theorem here that does not follow \\
\midrule
Overlapping decision regions~\cite{javdani2014} & The minimax law $\Theta(\log(1/\delta)/t^{*})$
for every correspondence and channel, and the order law $t^{*}=Ch^{2s^{*}}$
(Thms.~\ref{thm:lower}--\ref{thm:order}) \\
High-moment masking~\cite{grosso2014,duc2015} & The two minimax rates on one channel against
every almost-surely stopping policy, and uniformity over the linear class with a
query-free constant (Thms.~\ref{thm:endpoint}, \ref{thm:unif}) \\
Multiple-answer games~\cite{degenne2019,wangyou2026} & The obstruction-based lower bound, whose
hard pair is chosen from the policy's own output (Thm.~\ref{thm:lower}) \\
\bottomrule
\end{tabular}
\end{table}

\subsection{What is proved, and how}

Three evidence classes appear, and they are never interchanged.
\begin{itemize}
\item \PROVED: a theorem checked in Lean~4~\cite{demoura2021} against Mathlib~\cite{mathlib2020}.
\item \CERT: an exact rational certificate.
\item \NUM: a floating-point or Monte Carlo result.
\end{itemize}
Each theorem carries one line naming its declaration. The supplement gives the proof
architecture and every supporting declaration. All modules compile with no \texttt{sorry} and no
\texttt{native\_decide}, and their axioms lie within \{\texttt{propext}, \texttt{Classical.choice},
\texttt{Quot.sound}\}. A policy's \emph{cost} is its expected number of informative, that is
non-dummy, queries. A dummy query has the same law in every world. The cost equals
$\mathbb{E}\tau$ for a policy that issues no dummy before it stops, which every decoder here
satisfies.

\section{Safe-decision geometry and obstructions}
\label{sec:geometry}

Let $W$ be a finite set of worlds, $A$ a finite set of answers and
$\mathrm{safe}:W\to 2^{A}$ the safe answers. $H\subseteq W$ is an obstruction when no answer is
safe for all of $H$. Let $\mathrm{sep}(w,w')$ be the least order at which some admissible query
separates two distinct worlds. A set is \emph{together at order $s$} when no pair in it is
separated by order $s$.

\begin{theorem}[The safe-order law]
\label{thm:safelaw}
Order $s$ suffices, meaning every set together at order $s$ has a common safe answer, if and
only if every obstruction contains a pair separated by order $s$. Hence the least sufficient
order is
\[
s^{*}=\max_{H\ \text{obstruction}}\ \min_{w\ne w'\in H}\mathrm{sep}(w,w').
\]
When $\mathrm{sep}$ is read off a separation predicate on queries, the law is a statement about
queries.
\end{theorem}

\leanbox{\lean{safeAt\_iff}, \lean{safeAt\_mono}, \lean{safeAt\_zero\_of\_no\_unsafe} in
\lean{LeanPaperBSafeOrderV1}; \lean{together\_iff\_no\_separating\_query},
\lean{safeAt\_iff\_query}, \lean{sep\_unique} in \lean{LeanPaperBSepDerivedV1}}

Quantifying over all obstructions rather than minimal ones is not a weakening, because
obstructions are upward closed. If no obstruction exists, order $0$ suffices and the decision
costs nothing. Two cases matter below. For \emph{overlapping} safe sets,
$\mathrm{safe}(w)=A\setminus\{w\}$ with $A=W$, every two safe sets meet, the whole world set is
the only obstruction, and $s^{*}=\min\mathrm{sep}$. For \emph{identification},
$\mathrm{safe}(w)=\{w\}$, every pair is an obstruction and $s^{*}=\max\mathrm{sep}$.

\leanbox{\lean{coSafe\_pairwise\_meet}, \lean{coSafe\_univ\_unsafe}, \lean{coSafe\_proper\_safe},
\lean{coSafe\_safeAt\_iff} in \lean{LeanPaperBSafeOrderV1}; \lean{id\_order} in
\lean{LeanPaperBObstructionLawV1}}

The safe sets of the parity channel below are not posited. They are computed from the
compression loss. Retaining the singleton $\{a\}$ at world $w$ has loss $\frac{m-1}{2}$ exactly
when $a=w$, so the safe answers at tolerance $0$ are all $a\ne w$.

\leanbox{\lean{selectedRadius\_self}, \lean{selectedRadius\_other}, \lean{safe\_iff\_ne} in
\lean{LeanPaperBTipV1}}

\section{The obstruction law}
\label{sec:law}

Fix a channel: for each world and query a density, and one dummy query with the same law in
every world. Let $\kappa(w,w')\ge0$ bound the per-query divergence between two worlds. The
\emph{obstruction information} is
\[
t^{*}=\min_{H\ \text{obstruction}}\ \max_{w\ne w'\in H}\kappa(w,w').
\]
It exists whenever some obstruction exists and every world has a safe answer. It is the largest
threshold $t$ at which every obstruction contains a pair with $\kappa\ge t$. Some obstruction
has all its pairs at most $t^{*}$.

\leanbox{\lean{obstruction\_value}, \lean{resolves\_iff\_le}, \lean{two\_le\_card\_of\_unsafe} in
\lean{LeanPaperBObstructionLawV1}}

\begin{theorem}[The lower law]
\label{thm:lower}
Take any channel with $\mathrm{KL}_{q}(w\|w')\le\kappa(w,w')$ for every non-dummy query, and
any randomised adaptive policy with a finite horizon whose output is safe with probability at
least $1-\delta$ in every world. Then some world $w_{0}$ satisfies
\[
\mathrm{cost}(w_{0})\cdot t^{*}\ \ge\ \frac{1-\delta}{|W|-1}\log\frac1\delta-\log2 .
\]
The same holds with the per-query $\chi^{2}$ in place of $\kappa$.
\end{theorem}

\leanbox{\lean{obstruction\_pigeonhole}, \lean{info\_le\_cost}, \lean{obstruction\_lower},
\lean{obstruction\_lower\_star} in \lean{LeanPaperBObstructionLawV1};
\lean{obstruction\_lower\_chi} in \lean{LeanPaperBObstructionParityV1}}

\emph{Why.} Take an obstruction $H$ whose pairs all carry at most $t^{*}$, and a world
$w_{0}\in H$. Whenever the output is safe for $w_{0}$ it is unsafe for some other world of $H$,
because $H$ has no common safe answer. By pigeonhole, some $v\in H$ receives unsafe output with
probability at least $\frac{1-\delta}{|H|-1}$ at $w_{0}$ and at most $\delta$ at $v$. So the pair
is chosen from the policy's own output, not fixed in advance. An event of those two masses
forces the transcript divergence past $\frac{1-\delta}{|H|-1}\log\frac1\delta-\log2$. The
transcript divergence is at most the expected informative queries times the per-query bound.

\begin{theorem}[The tournament]
\label{thm:upper}
Let every obstruction contain a pair with $\kappa\ge t$. Test every pair with $\kappa\ge t$ by
an antisymmetric test, keep the worlds that win all their tests, and output a common safe
answer of the survivors.
\begin{itemize}
\item The survivors always have a common safe answer.
\item The output is safe at $w$ whenever $w$ wins its own tests.
\item Suppose each test on $n$ samples errs at world $w$ with probability at most
$e^{-cn\kappa(w,v)}$, and $n\ge\log(|W|/\delta)/(ct)$. Then the output is safe with probability
at least $1-\delta$ at every world, at a cost of at most $|W|^{2}n$ queries.
\end{itemize}
\end{theorem}

\leanbox{\lean{survivors\_untested}, \lean{decode\_safe}, \lean{decode\_error},
\lean{test\_budget}, \lean{tournament\_law} in \lean{LeanPaperBObstructionLawV1};
\lean{tournament\_satisfiable} in \lean{LeanPaperBObstructionParityV1}}

Two survivors never form a tested pair, since one of them lost. So the survivors contain no
tested pair, and every obstruction contains one, which means the survivors are not an
obstruction. The truth survives unless one of its own tests fails, and a union bound over those
tests gives the error. The test exponent is a property of a concrete decoder run. It is a
hypothesis of this theorem and is satisfiable. So the upper law is unconditional in its
combinatorics and conditional on the channel's pairwise testability.

\begin{theorem}[The order law]
\label{thm:order}
On a graded channel, $\kappa(w,w')=Ch^{2\,\mathrm{sep}(w,w')}$ with $0<h<1$. Every obstruction
contains a pair with $\kappa\ge Ch^{2s}$ if and only if order $s$ suffices. Hence
$t^{*}=Ch^{2s^{*}}$, and the minimax cost is $\Theta(\log(1/\delta)/h^{2s^{*}})$ whenever the
tests match the grading.
\end{theorem}

\leanbox{\lean{graded\_resolves\_iff}, \lean{graded\_tStar} in \lean{LeanPaperBObstructionLawV1};
\lean{graded\_lower\_chi} in \lean{LeanPaperBObstructionParityV1}}

\section{Both rates on one channel}
\label{sec:parity}

\paragraph{The channel} World $w\in\{1,\dots,m\}$ carries the mask $v(j,w)=-1$ if $j=w$ and $+1$
otherwise. Query $j$ returns
\[
Y_{j}=v(j,w)\prod_{\ell\le j}Z_{\ell}+\frac hj\sum_{\ell\le j}Z_{\ell}+\xi,
\]
with $Z$ uniform on the sign cube and latent, and $\xi\sim\mathcal N(0,\sigma^{2})$. The
characteristic functions of two worlds differ by $2\sin t\,(\sin(th/j))^{j}e^{-\sigma^{2}t^{2}/2}$
in magnitude. So query $j$ separates world $j$ at order $j$, and two worlds separate at
$\min(w,w')$.

\leanbox{\lean{parity\_cf\_difference}, \lean{parity\_gap\_order}, \lean{query\_separates\_iff},
\lean{parity\_s\_safe}, \lean{parity\_s\_id} in \lean{LeanPaperBParityV1};
\lean{channel\_sep\_is\_read\_off} in \lean{LeanPaperBSepDerivedV1}}

\begin{theorem}[The two rates]
\label{thm:endpoint}
Fix $m\ge3$ and $\sigma>0$, the tip loss at tolerance $0$, and the class of randomised adaptive
policies with a measurable stopping rule that stops almost surely, with measurable outputs. The
stopping time need not be bounded, nor its expectation finite. There are $0<c\le C$, $h_{0}>0$
and $\delta_{0}\in(0,1)$, depending only on $m$ and $\sigma$, such that for $0<h\le h_{0}$ and
$0<\delta\le\delta_{0}$
\begin{gather*}
c\,\frac{\log(1/\delta)}{h^{2}}\le N^{*}_{\rm safe}\le C\,\frac{\log(1/\delta)}{h^{2}},\\
c\,\frac{\log(1/\delta)}{h^{2(m-1)}}\le N^{*}_{\rm id}\le C\,\frac{\log(1/\delta)}{h^{2(m-1)}},
\end{gather*}
where $N^{*}$ is the minimax worst-world expected number of informative queries.
\end{theorem}

\leanbox{\lean{parity\_tip\_finish}, \lean{parity\_rates}, \lean{safeCert}, \lean{idCert} in
\lean{LeanPaperBParityFinishV1}; \lean{stopped\_parity\_rates}, \lean{stopped\_safe\_lower},
\lean{stopped\_id\_lower}, \lean{trunc\_correct}, \lean{costS\_ofPol} in
\lean{LeanPaperBStoppedParityV1}}

\begin{theorem}[The lower sides, from the law]
\label{thm:parlaw}
The parity channel is graded: every non-dummy query has
$\chi^{2}\le a_{m,\sigma}h^{2\,\mathrm{sep}(w,w')}$ with $\mathrm{sep}(w,w')=\min(w,w')$. Its least
safe order is $1$ for the overlapping answers and $m-1$ for identification. Theorem~\ref{thm:order}
therefore gives, for every finite-horizon policy that is safe with probability $1-\delta$, a
world with cost at least
$\bigl(\frac{1-\delta}{m-1}\log\frac1\delta-\log2\bigr)/(a_{m,\sigma}h^{2})$. For identification
the same holds with $h^{2(m-1)}$.
\end{theorem}

\leanbox{\lean{parity\_graded}, \lean{parity\_coSafe\_order}, \lean{parity\_id\_order},
\lean{parity\_safe\_from\_law}, \lean{parity\_id\_from\_law} in
\lean{LeanPaperBObstructionParityV1}}

\emph{The upper sides.} Two explicit decoders run bounded statistics, $\cos(t_{0}Y)$ or
$\sin(t_{0}Y)$ by the parity of $j$. Their per-test budgets come from an exponential
concentration bound proved on the channel. The safe decoder queries only order $1$. The
identification decoder queries every order up to $m-1$, and its cost is the sum over all its
tests. \emph{Unbounded stopping.} A stopping policy is truncated at a horizon by which every
world has stopped with probability at least $1-\delta$. The truncation is correct at level
$2\delta$ and costs no more. The finite-horizon lower bound then applies, and a finite-horizon
decoder is itself a stopping policy.

\leanbox{\lean{safe\_correct}, \lean{safe\_cost}, \lean{id\_correct}, \lean{id\_cost} in
\lean{LeanPaperBParityDecoderV1}}

\section{The query class decides the order}
\label{sec:family}

The parity channel's orders come from its queries. A second family shows the same law with the
query class itself as the lever. For $r\ge2$ and $u,v\in\{\pm1\}$,
\[
F_{u,v}(z)=3uz_{1}-2z_{2}-2z_{3}+v\prod_{\ell=4}^{r+3}z_{\ell}
\]
on $\{\pm1\}^{r+3}$, with the uniform reference and support
$K=\{(t,t,t,1,\dots,1)\}$. The losses on $K$, computed as half the oscillation of the retained
residual, make the safe sets at $u=\pm1$ disjoint for $\varepsilon<1$, while $v$ does not move
them. \CERT{} A query picks a linear probe $g(z)=a\cdot z$ with $\|a\|_{1}\le1$ and returns
$Y=F_{u,v}(Z)+hg(Z)+\xi$, with $Z$ latent. So $u$ decides the answer and $v$ is a nuisance that
identification must still resolve.

\begin{theorem}[Moment vanishing]
\label{thm:vanish}
Let $\chi$ be the top character on $\{\pm1\}^{n}$ and $g$ a polynomial of degree at most $d$.
Then $\mathbb{E}[\chi g^{j}]=0$ whenever $jd<n$.
\end{theorem}

\leanbox{\lean{moment\_vanishes\_below\_order}, \lean{top\_character\_orthogonal} in
\lean{LeanUnifiedQueryOrderV1}}

\begin{theorem}[Attainment]
\label{thm:attain}
For the averaged tail query, $\sum_{z}\bigl(\prod_{i}z_{i}\bigr)\bigl(\sum_{i}z_{i}\bigr)^{k}$
vanishes for $k<n$ and equals $2^{n}n!$ at $k=n$, so $\mathbb{E}[\chi_{n}g^{n}]=n!/n^{n}\ne0$. On
the family, the decision probe separates the two values of $u$ by $12h$ in the second moment,
and the tail probe separates the two values of $v$ by $2(r+1)\,r!/r^{r}\,h^{r}$ in the
$(r+1)$-st moment.
\end{theorem}

\leanbox{\lean{cubeSum\_below\_diag}, \lean{cubeSum\_diag}, \lean{tail\_query\_moment\_ne\_zero} in
\lean{LeanPaperBAttainmentV1}; \lean{decision\_moment\_gap}, \lean{identification\_moment\_gap}
in \lean{LeanPaperBRatesV1}}

\begin{theorem}[Per-step information, uniform over the linear class]
\label{thm:unif}
Consider every probe with shifts inside the box $\sum_{i}|a_{i}|\le1$, $\sum_{i}|ha_{i}|\le h_{0}$,
centres within $R$. Let $Q$ be the flat average, over the sign cube, of the probe's translates
of one base Gaussian mixture, and let $P$ weight the even patterns double and drop the odd
ones. Then $P-Q$ is exactly the $r$-fold difference, $P$ and $Q$ are densities up to one common
normaliser, and
\[
\int\frac{(P-Q)^{2}}{\ell}\ \le\ h^{2r}r^{-2r}C(\sigma,R+h_{0},r),
\]
with $C$ explicit and independent of the probe. So $\chi^{2}(P,Q)$ and $\mathrm{KL}(P\|Q)$ are
$O(h^{2r})$ uniformly over the class. The same bound, with a factor $4$, holds for pairs whose
difference is the $r$-fold difference of $\varphi(\cdot-1)-\varphi(\cdot+1)$.
\end{theorem}

\leanbox{\lean{iterDiff\_mixture\_envelope}, \lean{hellinger\_rate\_bound} in
\lean{LeanPaperBRatesV1}; \lean{worlds\_diff}, \lean{chi2\_uniform},
\lean{kl\_uniform\_over\_queries} in \lean{LeanPaperBWorldsV1}; \lean{kl\_le\_envelope} in
\lean{LeanPaperBKLV1}; \lean{pair\_chi2\_bound} in \lean{LeanPaperBPairRateV1}}

On the parity channel the box is discharged from the construction. The channel's queries
induce shifts with $\sum_{i}|a_{i}|=1$ exactly, and centres $\pm1$. On the latent family,
identifying its two $v$-worlds with such a pair, with the box discharged, is a modelling step
that is not done here.

\leanbox{\lean{queries\_inside\_box}, \lean{channel\_kl\_uniform} in \lean{LeanPaperBQueryClassV1}}

\emph{The class is the lever.} The restriction to the linear class is part of the theorem. At
$r=3$ the degree-$2$ probe $g=\frac12(z_{4}z_{5}+z_{6})$ has $g^{2}=\frac12(1+\chi)$. Squaring it
recreates the character a degree cap was meant to exclude, and
$\mathbb{E}_{+}Y^{3}-\mathbb{E}_{-}Y^{3}=3h^{2}$ exactly. \CERT{} A bounded moment contrast forces a
Hellinger separation, so $H^{2}\ge ch^{4}$ and no uniform $Ch^{6}$ bound holds on the
degree-capped class.

\leanbox{\lean{hellinger\_lower\_of\_moment\_contrast} in \lean{LeanPaperBHellingerV1}} One parameter of nonlinearity,
$g_{\eta}=(z_{4}+z_{5}+z_{6}+\eta z_{4}z_{5})/(3+\eta)$, opens a square channel. Its density
difference has an odd channel at $\eta h^{2}$ and an even one at $h^{3}$. The cross term
vanishes by parity, so the two scales add: $\eta^{2}h^{4}+h^{6}$. The crossover sits at
$\eta\asymp h$.

\leanbox{\lean{moment\_two}, \lean{moment\_three}, \lean{cross\_term\_vanishes},
\lean{two\_scale\_shape} in \lean{LeanPaperBCrossoverV1}}

\emph{What is open on this family.} The obstruction law applies to any channel written in
protocol form with finitely many queries. For this family the per-step bound is proved for the
constructed pair, but the family's own channel is neither identified with that pair nor
assembled in that form. So the adaptive lower bound is not instantiated here.
Its upper side is proved with Chebyshev's $1/\alpha$ rather than $\log(1/\delta)$. The two
exponents on this family are therefore per-step orders, not minimax rates.

\leanbox{\lean{budget\_lower}, \lean{safe\_upper}, \lean{sides\_do\_not\_match} in
\lean{LeanPaperBHeadlineV1}}

\section{Evidence}
\label{sec:evidence}

Every script self-checks against a proved statement and refuses to write output if a check
fails.

\begin{table}[t]
\caption{Evidence packages. Each script self-checks against a Lean statement.}
\label{tab:evidence}
\centering
\footnotesize
\resizebox{\columnwidth}{!}{%
\begin{tabular}{@{}llr@{}}
\toprule
Package & Script & Checks \\
\midrule
The two decoders, parity channel & \texttt{EXP\_paperB\_decoder} & 7/7 \\
\quad paired rerun, formalized fallback & \texttt{EXP\_paperB\_decoder\_leanfallback} & 10/10 \\
\quad figure from the formalized counts & \texttt{FIG\_paperB\_decoder\_formal} & 4/4 \\
Family, safe sets, moment identities & \texttt{EXP\_paperB\_construction} & 8/8 \\
Hellinger orders, real law & \texttt{EXP\_paperB\_hellinger\_real} & 6/6 \\
Per-step bound over the linear class & \texttt{EXP\_paperB\_uniform\_hellinger} & 6/6 \\
Envelope chain & \texttt{EXP\_paperB\_envelope\_rate} & 8/8 \\
Query-class crossover, $g_{\eta}$ & \texttt{EXP\_paperB\_crossover} & 8/8 \\
\bottomrule
\end{tabular}}
\end{table}

\begin{figure*}[t]
\centering
\includegraphics[width=0.92\textwidth]{paperB_decoder_formal.pdf}
\caption{The two decoders on the parity channel, from the formalized decoders' own counts.
The paired rerun used the formalized fallback with the same seed and draws, and no new sampling
was done. Left and centre: worst-world error against the \emph{total} budget, which for
identification sums over all $m-1$ tests, at $m=3$. The bounds are family-wise
Clopper--Pearson, each tail at $0.05/(2J)$ with $J=616$: upper (solid) and lower (dotted). A
budget counts as achieving $\delta=0.05$ only when its upper bound clears it. A budget is
bracketed from below only by a \emph{contiguous prefix} of tested budgets whose lower bounds all
exceed $\delta$, because a fixed decoder's error need not be monotone in $N$. Both the
certificates and the brackets are computed from the formalized decoders' counts. They coincide
with those of the script rule on all $16$ panels. Three $(m,h)$ cells certify both decoders,
all at $m=3$, with total budgets $25\to6400$, $100\to12800$ and $100\to51200$. This is an
achieved-cost separation and not an identification-rate curve. Right: the log--log slope of
the mean gap converges to the query index $j$ (exact). (\texttt{FIG\_paperB\_decoder\_formal},
4/4.)}
\label{fig:bdecoder}
\end{figure*}

\begin{figure}[t]
\centering
\includegraphics[width=\columnwidth]{paperB_envelope_rate.pdf}
\caption{The proved envelope chain on the latent family. Left: $\int\Delta^{2}/\ell$ against
$h$ for $r=2,3,4$, with measured orders $4.00$, $6.00$, $8.00$ and the proved bound dashed; the
gap is the explicit constant, not claimed sharp (\NUM). Right: the envelope ratio under its
dominating Gaussian.}
\label{fig:envrate}
\end{figure}

\begin{figure}[t]
\centering
\includegraphics[width=\columnwidth]{paperB_crossover.pdf}
\caption{The query class as the lever. Left: the measured order in $h$ against $\eta$, falling
from six to four. Right: $H^{2}$ against $h$ with the two-channel form
$c_{6}h^{6}+c_{4}\eta^{2}h^{4}$ fitted on the coarse half and drawn through the fine half (\NUM).
The additivity of the two channels is Theorem-level (cross term zero by parity); the fitted
constants are not.}
\label{fig:crossover}
\end{figure}

\section{Limitations}
\label{sec:limits}

\begin{enumerate}
\item[1)] \textbf{The upper law needs tests.} Theorem~\ref{thm:upper} assumes pairwise tests of
exponent $c\kappa$. On the parity channel the rates are attained by its own decoders, whose
concentration is proved. That a tournament's tests exist with that exponent on a general
channel is a property of the channel.
\item[2)] \textbf{Finite horizon in the general lower law.} Theorem~\ref{thm:lower} is stated
for finite-horizon policies. On the parity channel the extension to every almost-surely
stopping policy is proved by truncation. For a general channel the same truncation argument
applies but is not formalised.
\item[3)] \textbf{The latent family is a boundary.} There the per-step orders are proved
uniformly over the linear class, while the adaptive instantiation and an exponential-tail upper
bound are not. Its upper budget carries $1/\alpha$.
\item[4)] \textbf{Constants depend on $|W|$, $m$ and $r$.} The lower constant carries
$\frac{1}{|W|-1}$ and the tournament $|W|^{2}\log|W|$. Nothing is uniform as these grow.
\item[5)] \textbf{Finitely many worlds and queries}, throughout.
\end{enumerate}

\section{Conclusion}

A safe decision costs what its obstructions cost. Among the sets of worlds that share no safe
answer, find the one whose best-separated pair is worst. Its per-query information $t^{*}$
fixes the minimax cost at $\log(1/\delta)/t^{*}$, and on a graded channel it becomes an order,
$h^{2s^{*}}$. Identification is the case where every pair is an obstruction. On the parity
channel both rates are proved against every almost-surely stopping policy, and they separate
by a whole power of the signal, $h^{2}$ against $h^{2(m-1)}$. On a latent family the query
class itself sets the order.

\section*{Companion material}
The supplement gives the proof architecture of every theorem, the finite-alphabet transcript and
test results, the change of measure on filtrations, the envelope and the density identification
behind Theorem~\ref{thm:unif}, the geometry and density realisation of the orders, and the status
of results superseded during preparation. The claim map records each headline's model,
declaration, scope and evidence.

\begin{thebibliography}{00}

\bibitem{javdani2014} S. Javdani, Y. Chen, A. Karbasi, A. Krause, D. Bagnell, and
S. Srinivasa, ``Near optimal Bayesian active learning for decision making,'' in
\emph{Proc. 17th Int. Conf. Artificial Intelligence and Statistics (AISTATS)}, 2014,
pp.~430--438.

\bibitem{grosso2014} V. Grosso, F.-X. Standaert, and E. Prouff, ``Low entropy masking
schemes, revisited,'' in \emph{Smart Card Research and Advanced Applications (CARDIS)},
2013, pp.~33--43.

\bibitem{rioul2021} O. Rioul, W. Cheng, and S. Guilley, ``Cumulant expansion of mutual
information for quantifying leakage of a protected secret,'' in \emph{Proc. IEEE Int. Symp.
Information Theory (ISIT)}, 2021, pp.~2596--2601.

\bibitem{duc2015} A. Duc, S. Faust, and F.-X. Standaert, ``Making masking security
proofs concrete,'' IACR ePrint 2015/119, 2015.

\bibitem{poli2023} M. Poli et al., ``Ideal abstractions for decision-focused learning,''
arXiv:2303.17062, 2023.

\bibitem{degenne2019} R. Degenne and W. M. Koolen, ``Pure exploration with multiple correct
answers,'' in \emph{Advances in Neural Information Processing Systems 32}, 2019,
pp.~14591--14600.

\bibitem{wangyou2026} Q. Wang and W. You, ``Ranking-and-selection with multiple correct
answers and non-answerable estimates,'' arXiv:2606.21889, 2026.

\bibitem{mathlib2020} The mathlib Community, ``The Lean mathematical library,'' in
\emph{Proc. 9th ACM SIGPLAN Int. Conf. Certified Programs and Proofs (CPP)}, 2020,
pp.~367--381.

\bibitem{demoura2021} L. de Moura and S. Ullrich, ``The Lean 4 theorem prover and
programming language,'' in \emph{Proc. 28th Int. Conf. Automated Deduction (CADE)}, 2021,
pp.~625--635.

\end{thebibliography}

\end{document}
